\section{Base Station}
\label{sec:base-station}

The base station sits between detection on the drones
(Section~\ref{sec:detection}) and interception (Section~\ref{sec:intercept}).
It does not fly and has no camera of its own; it works on what the drones send
it. A drone that detects an object starts flying towards it at once and, during
its flight, sends image crops of the object and the direction to it every
$\Delta t$ seconds (Section~\ref{sec:target}). The base station sees the
reports of every drone at the same time, which no single drone does. This
section describes what the base station could do with them.

\textbf{Confirming the target.} The base station could run a stronger
classifier than the one on the drones, since it is not limited by the drones'
weight and power. Instead of one image, it could classify the sequence of
crops a drone sends, in which motion and shape over time help to tell a drone
from a bird, and combine the views of several drones that see the same
object. Since missing a hostile drone is worse than a detour, it would be
tuned so that it rarely misses a drone, and call a drone back only when it is
confident that the object is not a drone.

\textbf{Estimating the target's position.} One camera gives only the direction
to the target, not its distance. With the directions from two drones, the
base station could estimate the target's position by triangulation. Two
drones at positions $\mathbf{p}_1$ and $\mathbf{p}_2$ see the target along unit
directions $\mathbf{u}_1$ and $\mathbf{u}_2$; since the two lines rarely meet
exactly, one option is the midpoint of their closest points, the classic
midpoint method \cite{hartley1997triangulation, hartley2004multiple},
\begin{equation}
  \hat{\mathbf{x}} = \tfrac{1}{2}\left(\mathbf{p}_1 + s^\ast\mathbf{u}_1
                   + \mathbf{p}_2 + t^\ast\mathbf{u}_2\right),
  \quad
  (s^\ast, t^\ast) = \arg\min_{s,t}
  \left\|\mathbf{p}_1 + s\,\mathbf{u}_1 - \mathbf{p}_2 - t\,\mathbf{u}_2\right\|^2 .
  \label{eq:triangulation}
\end{equation}
A Kalman filter could smooth the estimates over time and give the target's
velocity as well, which the MPC in Section~\ref{sec:intercept} needs, with new
directions matched to the track by, for example, the nearest-neighbour rule
\cite{barshalom2001estimation}. With more drones seeing the target, the
estimate could use all of their directions.

\textbf{Choosing the drones.} The base station could decide which drones take
part in an intercept: confirming the drone that is already on its way, and
sending a second drone so that the target can be triangulated. It could pick
the drones that can reach the hostile drone's path earliest, for example with
a consensus-based auction between the drones \cite{choi2009consensus}.
